\section{دمج المشغلات بالنموذج الديناميكي}
\label{sec:actuator_dynamics}

إن النموذج الديناميكي المستنتج في المعادلة (\ref{eq:robot_matrix_form}) يصف العلاقة بين حركة الروبوت والعزوم الميكانيكية المطبقة على مفاصله. ومع ذلك، في التطبيقات الصناعية الفعلية، يتم تشغيل المفاصل باستخدام محركات كهربائية مقترنة بعلب سرعة (\LR{Gearboxes}). وبالتالي، فإن إشارة التحكم الحقيقية هي الجهد الكهربائي $(u)$ المطبق على المحركات. للحصول على تمثيل دقيق للنظام بالكامل، لا بد من دمج ديناميكية هذه المحركات مع ديناميكية الوصلات الميكانيكية.

\subsection{النموذج الرياضي لمحرك التيار المستمر (\LR{DC Motor Model})}
بإهمال التحريضية الكهربائية للمحرك (\LR{Armature Inductance}) لكون الثابت الزمني الكهربائي أصغر بكثير من الثابت الزمني الميكانيكي، يمكن كتابة معادلة العزم الصافي الذي يولده المحرك وينقله إلى مفصل الروبوت بالصيغة التالية:

\begin{equation}
\tau_{act} = \left(\frac{K N}{R_a}\right) u - \left(N^2 J_m\right) \ddot{q} - N^2 \left(f_m + \frac{K^2}{R_a}\right) \dot{q}
\label{eq:motor_torque}
\end{equation}

حيث:
\begin{itemize}
    \item $(u)$: جهد التحكم المطبق على المحرك.
    \item $(K)$: ثابت العزم والقوة المحركة الكهربائية العكسية (\LR{Back-EMF Constant}).
    \item $(N)$: نسبة التخفيض لعلبة السرعة (\LR{Gear Ratio}).
    \item $(R_a)$: المقاومة الكهربائية لملفات المحرك.
    \item $(J_m)$: عزم القصور الذاتي العائد لدوار المحرك.
    \item $(f_m)$: معامل الاحتكاك اللزج الخاص بالمحرك.
\end{itemize}

ولتبسيط التدوين الرياضي تمهيداً لدمج المعادلات، تم تجميع الثوابت الفيزيائية للمحرك في ثلاثة معاملات رئيسية $(A, B, c)$ كما يلي:
\begin{equation}
A = \frac{K N}{R_a} \quad , \quad B = N^2 J_m \quad , \quad c = N^2 \left(f_m + \frac{K^2}{R_a}\right)
\end{equation}
وبالتالي تصبح معادلة المحرك:
\begin{equation}
\tau_{act} = A u - B \ddot{q} - c \dot{q}
\label{eq:motor_simplified}
\end{equation}

\subsection{النموذج الديناميكي الشامل والمدمج (\LR{Comprehensive Dynamic Model})}
لإيجاد توازن العزوم الكلي، نقوم بمساواة العزم المحرك المولد من الآلية (\ref{eq:motor_simplified}) مع العزم المقاوم الذي يتطلبه الروبوت للحركة (\ref{eq:robot_matrix_form})، مع إضافة قوى الاحتكاك اللزج الميكانيكي للمفاصل الممثلة بالمعامل $(f)$. تصبح معادلة التوازن لكل مفصل $(i)$:

\begin{equation}
D_{i1}\ddot{q}_1 + D_{i2}\ddot{q}_2 + D_{i3}\ddot{q}_3 + f_i \dot{q}_i + D_i = A_i u_i - B_i \ddot{q}_i - c_i \dot{q}_i
\end{equation}

بنقل جميع الحدود المضروبة بالتسارع $(\ddot{q})$ والسرعة $(\dot{q})$ إلى طرف واحد، نلاحظ أن عطالة المحرك $(B)$ تُضاف مباشرة إلى القطر الرئيسي لمصفوفة القصور الذاتي للروبوت $(D_{ii})$، وأن احتكاك المحرك $(c)$ يُضاف إلى احتكاك المفصل $(f)$. لتسهيل الصياغة، نُعرّف المعاملات المدمجة التالية:
\begin{equation}
\bar{M}_{ii} = D_{ii} + B_i \quad , \quad R_i = f_i + c_i
\end{equation}

وبإعادة كتابة النظام بأكمله بصيغة شعاعية مصفوفية مدمجة، نحصل على:
\begin{equation}
\bar{M}(q) \begin{bmatrix} \ddot{\theta}_1 \\ \ddot{\theta}_2 \\ \ddot{d}_3 \end{bmatrix} = 
\begin{bmatrix} A_1 u_1 - R_1 \dot{\theta}_1 - D_1 \\ A_2 u_2 - R_2 \dot{\theta}_2 - D_2 \\ A_3 u_3 - R_3 \dot{d}_3 - D_3 \end{bmatrix}
\end{equation}

حيث تمثل $\bar{M}(q)$ مصفوفة العطالة الكلية للنظام (الروبوت + المحركات). وللحصول على شعاع التسارع كخرج صريح جاهز للاستخدام في محاكاة \LR{Simulink} وتصميم خوارزميات التحكم المتقدمة، نضرب طرفي المعادلة بمقلوب مصفوفة العطالة $\bar{M}^{-1}$:

\begin{equation}
\begin{bmatrix} \ddot{\theta}_1 \\ \ddot{\theta}_2 \\ \ddot{d}_3 \end{bmatrix} = \bar{M}^{-1} 
\begin{bmatrix} A_1 u_1 - R_1 \dot{\theta}_1 - D_1 \\ A_2 u_2 - R_2 \dot{\theta}_2 - D_2 \\ A_3 u_3 - R_3 \dot{d}_3 - D_3 \end{bmatrix}
\label{eq:final_state_space}
\end{equation}

تُعد المعادلة (\ref{eq:final_state_space}) تتويجاً لجهود النمذجة الديناميكية في هذا الفصل. فهذا النموذج الشامل، الذي يأخذ بعين الاعتبار تغيرات كتلة الحمل، وانزياح مراكز الثقل، والقصور الذاتي اللحظي، وديناميكية المحركات الكهربائية، يشكل بيئة الاختبار (\LR{Plant}) الدقيقة التي ستخضع لاستراتيجيات التحكم المتين والتكيفي في الفصول القادمة.